\documentclass[11pt]{article}
\usepackage[a4paper,margin=1in]{geometry}
\usepackage{amsmath,amssymb,amsthm,mathtools}
\usepackage{booktabs,longtable,array}
\usepackage{hyperref}
\usepackage{seqsplit}
\usepackage{xcolor}
\usepackage{enumitem}
\usepackage[T1]{fontenc}
\usepackage[utf8]{inputenc}

\hypersetup{colorlinks=true,linkcolor=blue,citecolor=blue,urlcolor=blue}
\newtheorem{theorem}{Theorem}
\newtheorem{lemma}{Lemma}
\newtheorem{proposition}{Proposition}
\newtheorem{corollary}{Corollary}
\newtheorem{remark}{Remark}
\newcommand{\Z}{\mathbb Z}
\newcommand{\bigint}[1]{\texttt{\seqsplit{#1}}}
\newcommand{\numblock}[2]{\par\noindent\( #1 \)\;\begin{minipage}[t]{0.82\linewidth}{\ttfamily\scriptsize\raggedright\seqsplit{#2}\par}\end{minipage}\par}
\newcommand{\modulo}[1]{\; (\mathrm{mod}\; #1)}
\setlist[itemize]{topsep=2pt,itemsep=2pt,parsep=0pt}
\setlist[enumerate]{topsep=2pt,itemsep=2pt,parsep=0pt}

\title{Congruence Conditions for the Square Constraint in\\
\texorpdfstring{$\displaystyle y^2=(6n+x)^2+\frac{36n^3-19}{x}$}{y2=(6n+x)2+(36n3-19)/x}}
\author{Prepared for Jamal Agbanwa}
\date{July 10, 2026}

\begin{document}
\maketitle

\begin{abstract}
Starting from the large CRT/cube-root congruences that force
\(x\mid 36n^3-19\), we analyze the extra condition that
\[
        y^2=(6n+x)^2+\frac{36n^3-19}{x}
\]
be an integer square.  The main conclusion is sharp but restrictive: the
large congruence construction must at least be refined to
\[
        n\equiv 0 \pmod 3,\qquad x\equiv 7 \pmod {12}.
\]
For the prime-product CRT construction with \(x=M=\prod_i p_i\), where every
\(p_i\equiv 2\pmod 3\), this means \(M\equiv 7\pmod {12}\); equivalently, the
number of prime factors \(p_i\equiv 5\pmod {12}\) and the number of prime
factors \(p_i\equiv 11\pmod {12}\) must both be odd.  The remaining square
condition is not a cosmetic congruence issue: it is exactly equivalent to a
representation of \(114\) as a sum of three integer cubes.  Thus the CRT
classes below are admissible search classes, not certificates of an integer
\(y\).
\end{abstract}

\section{Starting point: the divisibility congruence}
The uploaded note constructs large congruence classes by the following method.
For a prime \(p\equiv 2\pmod 3\), cubing is a bijection modulo \(p\). Hence
there is a unique residue \(a\pmod p\) satisfying
\[
        36a^3\equiv 19\pmod p.
\]
If \(p_1,\dots,p_k\) are distinct primes congruent to \(2\pmod 3\), and
\(a_i\) is the corresponding residue modulo \(p_i\), then the CRT gives a
unique residue \(A\pmod M\), with \(M=p_1\cdots p_k\), such that
\[
        n\equiv A\pmod M
        \quad\Longrightarrow\quad
        M\mid 36n^3-19.
\]
This document asks what must be added so that
\[
        (6n+M)^2+\frac{36n^3-19}{M}
\]
is also a square.

\section{Forced congruences from the square condition}
Let
\[
        q=\frac{36n^3-19}{x},\qquad z=6n+x.
\]
Then the equation is \(y^2=z^2+q\), and \(xq=36n^3-19\).

\begin{lemma}[Basic local restrictions]\label{lem:basic}
Suppose \(n,x,y\in\Z\), \(x\ne 0\), \(x\mid 36n^3-19\), and
\[
        y^2=(6n+x)^2+\frac{36n^3-19}{x}.
\]
Then
\[
        n\equiv 0\pmod 3,
        \qquad x\equiv 7\pmod {12},
        \qquad y\equiv 0\pmod 6.
\]
\end{lemma}

\begin{proof}
Since \(36n^3-19\) is odd and is congruent to \(-1\pmod 3\), every divisor
\(x\) is odd and prime to \(3\).  Thus \(q=(36n^3-19)/x\) is also odd and
prime to \(3\).

Modulo \(3\), we have \(6n+x\equiv x\) and \(36n^3-19\equiv -1\equiv 2\).
Therefore
\[
        y^2\equiv x^2+2x^{-1}\pmod 3.
\]
If \(x\equiv 2\pmod 3\), then \(x^2+2x^{-1}\equiv 1+2\cdot 2\equiv 2\pmod 3\),
which is not a square modulo \(3\). Hence \(x\equiv 1\pmod 3\), and then
\(y^2\equiv 0\pmod 3\), so \(3\mid y\).

Because \(3\mid y\), we have \(y^2\equiv 0\pmod 9\).  Also
\(36n^3-19\equiv -1\pmod 9\).  Reducing the equation modulo \(9\) and
multiplying by \(x\), we get
\[
        x(6n+x)^2\equiv 1\pmod 9.
\]
Since \(x\equiv 1\pmod 3\), the possible residues of \(x\pmod 9\) are
\(1,4,7\).  If \(n\equiv 0\pmod 3\), then \(6n\equiv 0\pmod 9\), and indeed
\(x(6n+x)^2\equiv x^3\equiv 1\pmod 9\). If \(n\equiv 1\pmod 3\), then
\(6n\equiv 6\pmod 9\), and for \(x=1,4,7\pmod 9\) one gets
\(x(x+6)^2\equiv 4\pmod 9\). If \(n\equiv 2\pmod 3\), then \(6n\equiv 3\pmod 9\),
and one gets \(x(x+3)^2\equiv 7\pmod 9\).  Therefore \(n\equiv 0\pmod 3\).

Modulo \(4\), since \(36n^3-19\equiv 1\pmod 4\) and \(x\) is odd, we have
\(q\equiv x^{-1}\equiv x\pmod 4\).  Also \(6n+x\) is odd, so
\((6n+x)^2\equiv 1\pmod 4\).  Thus
\[
        y^2\equiv 1+x\pmod 4.
\]
If \(x\equiv 1\pmod 4\), the right side is \(2\pmod 4\), impossible for a
square.  Hence \(x\equiv 3\pmod 4\), and then \(y^2\equiv 0\pmod 4\), so
\(2\mid y\).  Combining \(x\equiv 1\pmod 3\) and \(x\equiv 3\pmod 4\) gives
\(x\equiv 7\pmod {12}\), while \(3\mid y\) and \(2\mid y\) give \(6\mid y\).
\end{proof}

\section{Exact link with the sum of three cubes problem for 114}
The previous lemma gives unavoidable congruence filters.  The next theorem
explains why these filters do not finish the problem.

\begin{theorem}[Equivalence with three cubes]\label{thm:equiv}
There exist integers \(n,x,y\), with \(x\ne 0\), satisfying
\[
        y^2=(6n+x)^2+\frac{36n^3-19}{x}
        \quad\text{and}\quad x\mid 36n^3-19
\]
if and only if there exist integers \(U,V,W\) such that
\[
        U^3+V^3+W^3=114.
\]
More precisely, every solution \((n,x,y)\) gives a three-cube representation by
\[
        U=6n+2x,
        \qquad V=y-x,
        \qquad W=-y-x.
\]
Conversely, any representation \(U^3+V^3+W^3=114\) can be ordered so that the
even term is \(U\); then
\[
        n=\frac{U+V+W}{6},
        \qquad x=-\frac{V+W}{2},
        \qquad y=\frac{V-W}{2}
\]
gives an integer solution of the displayed equation.
\end{theorem}

\begin{proof}
Assume first that \((n,x,y)\) is a solution.  Put
\[
        U=6n+2x,
        \quad V=y-x,
        \quad W=-y-x.
\]
Then
\[
\begin{aligned}
U^3+V^3+W^3
&=(6n+2x)^3+(y-x)^3+(-y-x)^3\\
&=(6n+2x)^3-6xy^2-2x^3\\
&=6\bigl(36n^3+36n^2x+12nx^2+x^3-xy^2\bigr).
\end{aligned}
\]
The original equation is equivalent, after multiplying by \(x\), to
\[
        xy^2=36n^3+36n^2x+12nx^2+x^3-19.
\]
Substitution into the previous line gives
\[
        U^3+V^3+W^3=6\cdot 19=114.
\]

Conversely, suppose \(U^3+V^3+W^3=114\).  Since cubes modulo \(9\) are
\(0,\pm 1\), and \(114\equiv -3\pmod 9\), each of \(U,V,W\) must have cube
\(-1\pmod 9\), so each is congruent to \(2\pmod 3\).  Also, since \(114\) is
even but not divisible by \(8\), exactly one of \(U,V,W\) is even.  Reorder the
three variables so that \(U\) is the even one.  Then
\[
        U\equiv 2\pmod 6,
        \qquad V\equiv W\equiv 5\pmod 6,
\]
so \(U+V+W\equiv 0\pmod 6\), and \(V+W\), \(V-W\) are even.  Therefore
\[
        n=\frac{U+V+W}{6},\quad
        x=-\frac{V+W}{2},\quad
        y=\frac{V-W}{2}
\]
are integers.  The identity computed above reverses to give
\[
        xy^2=36n^3+36n^2x+12nx^2+x^3-19
              =x(6n+x)^2+36n^3-19.
\]
Thus
\[
        x\bigl(y^2-(6n+x)^2\bigr)=36n^3-19.
\]
The case \(x=0\) would force \(V=-W\) and hence \(U^3=114\), impossible in
integers.  Thus \(x\ne 0\), and the last identity gives both divisibility by
\(x\) and the required equation.
\end{proof}

\begin{corollary}\label{cor:notcongruenceonly}
No finite congruence refinement of the CRT divisibility construction should be
treated as a certificate that \(y\) is integral.  A successful instance is
equivalent to finding an integer representation of \(114\) as a sum of three
cubes.
\end{corollary}

\section{How to refine the large CRT construction}
The uploaded CRT construction already gives divisibility.  Lemma~\ref{lem:basic}
shows that any candidate class must also satisfy
\[
        n\equiv 0\pmod 3,
        \qquad x\equiv 7\pmod {12}.
\]
For the prime-product construction, take
\[
        M=\prod_{i=1}^k p_i,
        \qquad p_i\equiv 2\pmod 3,
\]
and use \(x=M\).  Since each odd prime \(p_i\equiv 2\pmod 3\) is congruent to
\(5\) or \(11\pmod {12}\), the condition \(M\equiv 7\pmod {12}\) is equivalent
to
\[
        \#\{i:p_i\equiv 5\pmod {12}\}\equiv 1\pmod 2,
        \qquad
        \#\{i:p_i\equiv 11\pmod {12}\}\equiv 1\pmod 2.
\]
The smallest such choice is one prime of each type.

Thus an admissible search class has the form
\[
\boxed{
\begin{aligned}
        n&\equiv A \pmod {3M},\\
        x&=M,\qquad M\equiv 7\pmod {12},
\end{aligned}}
\]
where \(A\) is the CRT solution of
\[
        A\equiv 0\pmod 3,
        \qquad A\equiv a_i\pmod {p_i}\quad(1\le i\le k),
\]
and each \(a_i\) satisfies \(36a_i^3\equiv 19\pmod {p_i}\).
For every \(n\equiv A\pmod {3M}\), one has \(M\mid 36n^3-19\), and the forced
modulo \(3,4,9,12\) obstructions to \(y\) being an integer square are removed.
The remaining condition is the genuine three-cubes condition from
Theorem~\ref{thm:equiv}.


\section{Explicit admissible large congruence skeletons}
The following examples use prime/root data from the uploaded large-congruence
computations, but select products \(M\equiv 7\pmod {12}\) and add the necessary
condition \(n\equiv 0\pmod 3\).

\subsection*{A 60-digit admissible divisor}
Use the following two prime/root pairs:
\numblock{p_1=}{779282871203804318581334453057}
\numblock{a_1=}{540863547358444129637446662643}
\numblock{p_2=}{422843868764454355911216773723}
\numblock{a_2=}{173885648024960625137792429386}
Here \(p_1\equiv 5\pmod {12}\), \(p_2\equiv 11\pmod {12}\), and both are
\(2\pmod 3\).  The product is \(M_{60}=p_1p_2\), where
\numblock{M_{60}=}{329514984121688619762783660355521561208464753917624334621211}
so \(M_{60}\equiv 7\pmod {12}\).  The refined CRT class is
\[
        n\equiv A_{60}\pmod {N_{60}},
        \qquad N_{60}=3M_{60},
\]
where
\numblock{A_{60}=}{6123585947301432250420495038918585980078995990135335295170}
\numblock{N_{60}=}{988544952365065859288350981066564683625394261752873003863633}
For every \(n\) in this class and \(x=M_{60}\), one has \(x\mid 36n^3-19\),
while \(n\equiv 0\pmod 3\) and \(x\equiv 7\pmod {12}\).

\subsection*{An 80-digit admissible divisor}
Use the following two prime/root pairs:
\numblock{p_1=}{7729484335457653901640057298531371241781}
\numblock{a_1=}{7668575607239450973459863267707132263860}
\numblock{p_2=}{4350736286376564530047381650448847768399}
\numblock{a_2=}{525176878515189528938215354529508172272}
Here \(p_1\equiv 5\pmod {12}\), \(p_2\equiv 11\pmod {12}\), and both are
\(2\pmod 3\).  The product is \(M_{80}=p_1p_2\), where
\numblock{M_{80}=}{33628947973254860882585081701145824718581850550296972738720431240616807320278619}
so \(M_{80}\equiv 7\pmod {12}\).  The refined CRT class is
\[
        n\equiv A_{80}\pmod {N_{80}},
        \qquad N_{80}=3M_{80},
\]
where
\numblock{A_{80}=}{95585981323865432505725289557587562987151849436683367780975283125799528610477321}
\numblock{N_{80}=}{100886843919764582647755245103437474155745551650890918216161293721850421960835857}
For every \(n\) in this class and \(x=M_{80}\), the quotient \((36n^3-19)/x\) is
an integer and the necessary square congruences from Lemma~\ref{lem:basic} are
satisfied.

For the least representative \(n=A_{80}\) in this 80-digit class, exact
arithmetic gives
\numblock{(36A_{80}^3-19)/M_{80}=}{934914356623755878709556229742361127659491904728230197026247469846131050286142001010783580690572385245334842196936454823101781996033924286298169758151340324379683}
but
\[
        (6A_{80}+M_{80})^2+\frac{36A_{80}^3-19}{M_{80}}
\]
is not a square.  This is not a contradiction; it means only that the least
representative of the admissible class is not a solution.

\section{Interpretation}
The answer to the congruence question is therefore:
\[
        \boxed{n\equiv 0\pmod 3,\\ x\equiv 7\pmod {12}}
\]
must be imposed on top of the large divisibility congruences.  In the CRT
prime-product method, choose \(M\equiv 7\pmod {12}\), extend the CRT class for
\(n\) by adding \(n\equiv 0\pmod 3\), and then search within
\(n=A+3Mt\), \(x=M\), for a value of \(t\) making the displayed quantity a
square.

However, by Theorem~\ref{thm:equiv}, a fully successful integer triple is
nothing less than a representation of \(114\) as a sum of three integer cubes.
The congruences above identify admissible lanes for a search; they do not by
themselves solve the square condition.

\section*{References}
\begin{enumerate}
\item \emph{Large Modular Congruences for which \(x\mid 36n^3-19\): A CRT / cube-root construction with 30- and 40-digit moduli}, uploaded note, July 10, 2026.
\item Andrew V. Sutherland, \emph{Sums of three cubes}, Waterloo slides, 2019.
\item Public status summaries for the sums of three cubes problem list \(114\) among the remaining unsolved cases up to \(1000\) as of the sources checked on July 10, 2026.
\end{enumerate}

\end{document}
